\exercisesection{分解}

\begin{enumerate}
\def\labelenumi{\arabic{enumi})}
\tightlist
\item
  设 \(a, b\) 为 A 的两个元素，使得 \(a\) （相应地， \(b\) ）的左零化子是理想 \(Ab\) （相应地， \(Aa\) ）。证明序列
\end{enumerate}

\[
\dots \longrightarrow \mathrm{A} _ {s} \xrightarrow {\delta_ {a}} \mathrm{A} _ {s} \xrightarrow {\delta_ {b}} \mathrm{A} _ {s} \xrightarrow {\delta_ {a}} \mathrm{A} _ {s} \longrightarrow \dots \longrightarrow \mathrm{A} _ {s} \xrightarrow {\delta_ {a}} \mathrm{A} _ {s} \longrightarrow \mathrm{A} / \mathrm{A} a \longrightarrow 0,
\]

其中 \(\delta_{a}\) （相应地， \(\delta_{b}\) ）表示右乘 \(a\) （相应地， \(b\) ），定义了 A-模 A/Aa 的一个自由分解。证明此构造适用于以下两种情形：

\(\alpha)\) 设 \(\mathbf{A}_0\) 为一个环；我们取 \(\mathrm{A} = \mathrm{A}_0(\varepsilon)\) （X，第 27 页）， \(a = b = \varepsilon\) 。

β）设G为有限循环群，σ为G的一个生成元，k为交换环。我们取

\[
\mathbf {A} = k ^ {(\mathbf {G})}, \quad a = 1 - e _ {\sigma}, \quad b = \sum_ {g \in \mathbf {G}} e _ {g}.
\]

\begin{enumerate}
\def\labelenumi{\arabic{enumi})}
\setcounter{enumi}{1}
\tightlist
\item
  我们假设环 A 是交换的。设 B 为一个 A-代数；对 \(n \geqslant 0\) ，我们用 \(\mathcal{A}^{n}(\mathbf{B})\) 表示 \((n + 1)\) 个都等于 B 的模在 A 上的张量积。我们设 \(\mathcal{A}^{n}(\mathbf{B}) = 0\) （当 \(n < 0\) ），并设 \(\mathcal{A}(\mathbf{B}) = \bigoplus \mathcal{A}^{n}(\mathbf{B})\) 。我们通过如下设定，定义 \(\mathcal{A}(\bar{\mathbf{B}})\) 的一个次数递增 +1 的分次 A-自同态：
\end{enumerate}

\[
d (b _ {0} \otimes \dots \otimes b _ {n}) = \sum_ {i = 0} ^ {n + 1} (- 1) ^ {i} b _ {0} \otimes \dots \otimes b _ {i - 1} \otimes 1 _ {\mathbf {B}} \otimes b _ {i} \otimes \dots \otimes b _ {n} \quad \text { for } b _ {0}, \dots , b _ {n} \in \mathbf {B}.
\]

\begin{enumerate}
\def\labelenumi{\alph{enumi})}
\tightlist
\item
  证明 \(d \circ d = 0\) ，使得 \((\mathcal{A}(\mathbf{B}), d)\) 是 A-模的一个复形。若 M 是 A-模，我们记 \(\mathcal{A}(\mathbf{B}, \mathbf{M})\) 为满足 \(\mathcal{A}^{n}(\mathbf{B}, \mathbf{M}) = \mathcal{A}^{n}(\mathbf{B}) \otimes_{\mathbf{A}} \mathbf{M}\) 的复形，其微分为 \(d \otimes 1_{\mathbf{M}}\) 。 b) 假设代数 B 是忠实平坦的 A-模（X，第 169 页，习题 9）。证明复形 \(\mathcal{A}(\mathbf{B}, \mathbf{M})\) 定义了 A-模 M 的一个右分解。（只需验证复形 \(\mathbf{B} \otimes_{\mathbf{A}} \mathcal{A}(\mathbf{B}, \mathbf{M})\) 定义了 \(\mathbf{B} \otimes_{\mathbf{A}} \mathbf{M}\) 的一个分解；构造这两个复形之间的一个 A-线性同伦等价。）
\end{enumerate}

\begin{enumerate}
\def\labelenumi{\arabic{enumi})}
\setcounter{enumi}{2}
\tightlist
\item
  设 B 是一个环，a 是 B 的一个双边理想，b 是 B 的一个包含 a 的左理想。记 A 为环 B/a，用 \(p: \mathbf{B} \to \mathbf{A}\) 表示商同态，用 \(b'\) 表示理想 \(p(\mathbf{b}) \subset \mathbf{A}\) 。a) 考虑典范嵌入的如下序列：
\end{enumerate}

\[
\dots \rightarrow a ^ {n} b \rightarrow a ^ {n} \rightarrow a ^ {n - 1} b \rightarrow a ^ {n - 1} \rightarrow \dots \rightarrow a \rightarrow b \rightarrow B.
\]

证明通过与B/a作张量积得到的序列

\[
\dots \rightarrow \mathfrak {a} ^ {n} \mathfrak {b} / \mathfrak {a} ^ {n + 1} \mathfrak {b} \rightarrow \mathfrak {a} ^ {n} / \mathfrak {a} ^ {n + 1} \rightarrow \dots \rightarrow \mathfrak {a} / \mathfrak {a} ^ {2} \rightarrow \mathfrak {b} / \mathfrak {a b} \rightarrow \mathbf {A}
\]

定义了A-模A/b'的一个左分解。

\begin{enumerate}
\def\labelenumi{\alph{enumi})}
\setcounter{enumi}{1}
\item
  假设理想 a 和 b 是自由左 A-模。证明前述分解是 A-模 A/b' 的一个自由分解。
\item
  设 G 为一个群，k 为一个交换环， \(A = \dot{k}^{(G)}\) 为 G 在 k 上的代数， \(\varepsilon : A \to k\) 为 k-代数的同态，使得 \(\varepsilon(e_g) = 1\) 对所有 \(g \in G\) 成立。我们赋予 k 与同态 \(\varepsilon\) 相伴的 A-模结构（左模结构）。设 \((g_i)_{i \in I}\) 为 G 的一个生成族，F 为在 I 上构造的自由群， \(B = k^{(F)}\) , \(\pi : F \to G\) 为满足 \(\pi(i) = g_i\) （当 \(i \in I\) ）的同态， \(p : B \to A\) 为由 \(\pi\) 导出的 k-代数同态。证明理想 a = Ker(p) 与 b = Ker ( \(\varepsilon \circ p\) ) 是自由左 B-模；由此推导出 \(k^{(G)}\) -模 k 的一个自由分解。
\end{enumerate}

（根据 I，第 147 页，习题 20，群 Ker (π) 具有一个基本族 (rα)α∈J；证明元素 (eᵣα - 1)，其中 α ∈ J（相应地 (eᵢ - 1)，其中 i ∈ I）构成左 B-模 a（相应地 b）的一组基。）

\begin{enumerate}
\def\labelenumi{\alph{enumi})}
\setcounter{enumi}{3}
\tightlist
\item
  采用 c) 中的记号，我们设 \(\mathbf{R} = \operatorname{Ker}(\pi)\) ；令 \(\varphi: \mathbf{R} \to \mathbf{R}/(\mathbf{R}, \mathbf{R})\) 为商映射。我们赋予 \(k\) -模 \(k \otimes_{\mathbf{Z}} \mathbf{R}/(\mathbf{R}, \mathbf{R})\) 一个 A-模结构，使得
\end{enumerate}

\[
e _ {\pi (f)} \left(1 \otimes \varphi (r)\right) = 1 \otimes \varphi (f r f ^ {- 1})
\]

对所有 \(r \in \mathbb{R}\) , \(f \in \mathbb{F}\) 成立。证明存在 A-模的一个正合列：

\[
0 \rightarrow k \otimes_ {\mathbf {Z}} \mathrm{R} / (\mathrm{R}, \mathrm{R}) \rightarrow \mathrm{A} ^ {(I)} \rightarrow \mathrm{A} \rightarrow k \rightarrow 0.
\]

\begin{enumerate}
\def\labelenumi{\arabic{enumi})}
\setcounter{enumi}{3}
\tightlist
\item
  假设环 A 是交换的。
\end{enumerate}

设 G 为一个群，K 为一个单纯复形（X，第 177 页，习题 19）。G 在 K 上的作用是从 G 到由 K 到自身的双射单纯映射组成的群的一个同态。a) G 在 K 上的作用定义了 G 在 A-模 S(K, A) 上的作用（同上）；证明该作用使 S(K, A) 成为 A \(^{(G)}\) -模的一个复形。

\begin{enumerate}
\def\labelenumi{\alph{enumi})}
\setcounter{enumi}{1}
\item
  假设 G 在 K 的底集合上自由作用。证明 A \(^{(G)}\) -模 S \(_{n}\) (K, A) 是自由的。
\item
  假设单纯复形 K 是锥形的。证明 \(\mathbf{A}^{(\mathbf{G})}\) -复形 \(\mathbf{S}(\mathbf{K},\mathbf{A})\) 定义了 A 的一个左分解，其中 A 配备了 \(\mathbf{A}^{(\mathbf{G})}\) -模结构，使得 \(e_g a = a\) 对所有 \(g \in \mathbf{G}\) , \(a \in \mathbf{A}\) 成立。
\item
  设 \(|\mathbf{G}|\) 为单纯复形 \((\mathbf{G},\mathcal{F})\) ，其中 \(\mathcal{F}\) 是 G 的有限子集组成的集合。群 G 通过左平移作用在单纯复形 \(|\mathbf{G}|\) 上；证明 \(\mathbf{A}^{(\mathbf{G})}\) -复形 S(\textbar G\textbar, A) 是 \(\mathbf{A}^{(\mathbf{G})}\) -模 A 的一个自由分解。
\item
  设 \(\mathbf{B}(\mathbf{A}^{(G)},\mathbf{A})\) 为 \(\mathbf{A}^{(G)}\) -模 A 的标准分解；对 \(n\geqslant 0\) ， \(\mathbf{A}^{(G)}\) -模 \(\mathbf{B}_n(\mathbf{A}^{(G)},\mathbf{A})\) 等同于 \((\mathbf{A}^{(G)})^{\otimes (n + 1)}\) 。证明 A-线性同态 \(u:\mathrm{S}(|\mathrm{G}|,\mathrm{A})\to \mathrm{B}(\mathrm{A}^{(G)},\mathrm{A})\) （定义为
\end{enumerate}

\[
u \left(g _ {0}, \dots , g _ {n}\right) = g _ {0} \otimes g _ {0} ^ {- 1} g _ {1} \otimes \dots \otimes g _ {n - 1} ^ {- 1} g _ {n} \quad \text { for } g _ {0}, \dots , g _ {n} \text { in } G,
\]

）是 \(A^{(G)}\) -复形之间的一个同构。

\begin{enumerate}
\def\labelenumi{\arabic{enumi})}
\setcounter{enumi}{4}
\tightlist
\item
  \begin{enumerate}
  \def\labelenumii{\alph{enumii})}
  \tightlist
  \item
    设 \(0 \to K \to P \to M \to 0\) , \(0 \to K' \to P' \to M \to 0\) 为两个 A-模的正合列，其中 \(P\) 与 \(P'\) 是投射的。证明存在一个交换图：
  \end{enumerate}
\end{enumerate}

\[
\begin{array}{c c c}0 \rightarrow \mathrm{K} \oplus \mathrm{P} ^ {\prime} \rightarrow \mathrm{P} \oplus \mathrm{P} ^ {\prime} \rightarrow \mathrm{M} \rightarrow 0\\\downarrow^ {v}&\downarrow^ {u}&\downarrow^ {l _ {M}}\\0 \rightarrow \mathrm{P} \oplus \mathrm{K} ^ {\prime} \rightarrow \mathrm{P} \oplus \mathrm{P} ^ {\prime} \rightarrow \mathrm{M} \rightarrow 0\end{array}
\]

其中水平行是正合的，并且 \(u\) 与 \(v\) 是同构（参见 II，第 183 页，习题 4）。b) 设 \(n\) 为一个整数 \(\geqslant 0\) , \((\mathbf{P}, p)\) 与 \((\mathbf{P}', p')\) 为 A-模 M 的两个左分解，使得 \(\mathbf{P}_j = \mathbf{P}_j' = 0\) （当 \(j > n\) ）且 A-模 \(\mathbf{P}_i\) 与 \(\mathbf{P}_i'\) 对 \(0 \leqslant i \leqslant n - 1\) 是投射的。证明 A-模 \(\bigoplus_{i \geqslant 0} (\mathbf{P}_{2i} \oplus \mathbf{P}_{2i+1}')\) 与 \(\bigoplus_{i \geqslant 0} (\mathbf{P}_{2j}^{\prime} \oplus \mathbf{P}_{2j+1})\) 是同构的。

\begin{enumerate}
\def\labelenumi{\alph{enumi})}
\setcounter{enumi}{2}
\tightlist
\item
  在 \(b\) ) 的假设下，证明存在两个零同调的投射复形 \(\mathbf{R}\) , \(\mathbf{R}'\) ，使得 \(\mathbf{R}_i = \mathbf{R}_i' = 0\) （当 \(i \notin [0, n]\) ），以及分解之间的一个同构 \(u: \mathbf{P} \oplus \mathbf{R} \to \mathbf{P}' \oplus \mathbf{R}'\) 。
\end{enumerate}

¶6). 设 M 为一个左 A-模。称长度为 n 的表示，或 M 的 n-表示，为一个正合列

\[
\mathrm{L} _ {n} \rightarrow \mathrm{L} _ {n - 1} \rightarrow \dots \rightarrow \mathrm{L} _ {1} \rightarrow \mathrm{L} _ {0} \rightarrow \mathrm{M} \rightarrow 0
\]

其中 \(L_{i}\) 是自由左 A-模（ \(0 \leqslant i \leqslant n\) ）。称该表示为有限的，如果所有 \(L_{i}\) 都是有限生成的自由模。

如果 M 是一个有限生成的左 A-模，则用 \(\lambda(M)\) （有限或等于 \(+\infty\) ）表示整数 \(n \geqslant 0\) （M 对之具有有限 \(n\) -表示）的上确界。如果 M 不是有限生成的，则令 \(\lambda(M) = -1\) 。

\begin{enumerate}
\def\labelenumi{\alph{enumi})}
\item
  设 \(0 \to \mathbf{P} \to \mathbf{N} \to \mathbf{M} \to 0\) 为左 A-模的一个正合列。则有 \(\lambda(\mathbf{N}) \geqslant \inf (\lambda(\mathbf{P}), \lambda(\mathbf{M}))\) 。
\item
  设 \(L_{n} \xrightarrow{d_{n}} L_{n-1} \to \cdots \to L_{0} \to M \to 0\) 为 M 的一个有限 n-表示；如果 \(\lambda(M) > n\) ，证明 Ker \((d_{n})\) 是有限生成的 A-模（利用习题 5，b））。
\item
  证明：在 a) 的假设下，有
\end{enumerate}

\[
\lambda (\mathbf {M}) \geqslant \inf \left(\lambda (\mathbf {N}), \lambda (\mathbf {P}) + 1\right).
\]

\((\mathrm{If} n \leqslant \inf (\lambda(\mathbf{N}), \lambda(\mathbf{P}) + 1)\) ，通过对 \(n\) 作归纳证明 \(\lambda(\mathbf{M}) \geqslant n\) ，推理方式如同在 \(a\) 中并使用 \(b)\) 。）

\begin{enumerate}
\def\labelenumi{\alph{enumi})}
\setcounter{enumi}{3}
\tightlist
\item
  证明：在 a) 的假设下，有
\end{enumerate}

\[
\lambda (\mathbf {P}) \geqslant \inf \left(\lambda (\mathbf {N}), \lambda (\mathbf {M}) - 1\right).
\]

（方法与 c）相同。）由此推出，若 \(\lambda(\mathbf{N}) = +\infty\) ，则 \(\lambda(\mathbf{M}) = \lambda(\mathbf{P}) + 1\) 。

\begin{enumerate}
\def\labelenumi{\alph{enumi})}
\setcounter{enumi}{4}
\tightlist
\item
  由 a)、c) 和 d) 推出，若 \(\mathbf{N} = \mathbf{M} \oplus \mathbf{P}\) ，则有
\end{enumerate}

\[
\lambda (\mathbf {N}) = \inf \left(\lambda (\mathbf {M}), \lambda (\mathbf {P})\right).
\]

特别地，为使 N 具有有限表示，必要且充分的是 M 与 P 也具有有限表示。f) 设 \(N_{1}\) , \(N_{2}\) 为 A-模 M 的两个子模。假设 \(N_{1}\) 与 \(N_{2}\) 都具有有限表示。为使 \(N_{1} + N_{2}\) 具有有限表示，必要且充分条件是 \(N_{1} \cap N_{2}\) 是有限生成的。

\begin{enumerate}
\def\labelenumi{\arabic{enumi})}
\setcounter{enumi}{6}
\tightlist
\item
  \begin{enumerate}
  \def\labelenumii{\alph{enumii})}
  \tightlist
  \item
    使用习题 6 的记号，证明：若 M 是投射模，则有
  \end{enumerate}
\end{enumerate}

\[
\lambda (\mathbf {M}) = - 1 \quad \text { or } \quad \lambda (\mathbf {M}) = + \infty .
\]

若 A 是左诺特环，则对每个 A-模 M，有 \(\lambda(M) = -1\) 或 \(\lambda(M) = +\infty\) 。 b) 若 a 是环 A 的一个左理想，且不是有限生成的，则 \(A_s/a\) 是一个单演 A-模，它不具有有限表示，换句话说 \(\lambda(A_s/a) = 0\) （参见 X，第 11 页，命题 6）。

\begin{enumerate}
\def\labelenumi{\alph{enumi})}
\setcounter{enumi}{2}
\item
  给出环 A 的一个单演左理想 a 的例子，使得 \(A_{s}/a\) （它是有限表示的）容许一个对偶，该对偶不是有限生成的右 A-模。
\item
  设 K 为一个交换域，E 为向量空间 \(\mathbf{K}^{(\mathbf{N})}\) , \((e_n)\) 为 E 的典范基，T 为 E 的张量代数，因此它的一组基由有限乘积 \(e_{i_1} e_{i_2} \ldots e_{i_k}\) 构成（ \(k \geqslant 0\) , \(i_j \in \mathbf{N}\) 对所有 \(j\) ）。对给定的整数 \(n\) ，设 b 为 T 的由乘积 \(e_1 e_0, e_2 e_1, \ldots, e_n e_{n-1}\) 与 \(e_{n+k} e_n\) （对所有 \(k \geqslant 1\) ）生成的双边理想；设 A 为商环 T/b，并对每个整数 \(m\) ，设 \(a_{m}\) 为 \(e_{m}\) 在 A 中的典范像。证明，若 \(\mathbf{M} = \mathbf{A}_{s}/\mathbf{A}a_{0}\) ，则 \(\lambda(\mathbf{M}) = n\) （注意，对 \(m \leqslant n - 1\) ， \(a_{m}\) 的左零化子是 \(\mathbf{A}a_{m+1}\) ，并利用习题 6，b)）。
\end{enumerate}

\begin{enumerate}
\def\labelenumi{\arabic{enumi})}
\setcounter{enumi}{7}
\tightlist
\item
  设 C 为一个交换环，E、F 为两个 C-模。证明有 \(\lambda(E \otimes_{C} F) \geqslant \inf(\lambda(E), \lambda(F))\) 。
\item
  设 \(\rho: A \to B\) 为一个环同态，M 是一个 A-模。
\end{enumerate}

\begin{enumerate}
\def\labelenumi{\alph{enumi})}
\item
  如果 B 是平坦的右 A-模且 M 具有有限 n-表示，则 B-模 B \(\otimes_{\mathbf{A}}\) M 具有有限 n-表示。
\item
  如果 B 是忠实平坦的右 A-模（X，第 169 页，习题 9），并且 B-模 B ⊗\_A M 具有有限 n-表示，则 M 具有有限 n-表示（利用习题 6，b）和第 169 页习题 12）。
\end{enumerate}

\begin{enumerate}
\def\labelenumi{\arabic{enumi})}
\setcounter{enumi}{9}
\tightlist
\item
  设 E 是一个 A-模。称 E 是伪凝聚的，如果 E 的每个有限生成子模都是有限表示的；伪凝聚模的每个子模都是伪凝聚的。称 E 是凝聚的，如果它是伪凝聚的且有限生成的（因此是有限表示的）。
\end{enumerate}

\begin{enumerate}
\def\labelenumi{\alph{enumi})}
\item
  设 \(0 \to E' \to E \to E'' \to 0\) 为 A-模的一个正合列。证明，如果 E 是伪凝聚的（相应地，凝聚的）且 \(E'\) 是有限生成的，则 \(E''\) 是伪凝聚的（相应地，凝聚的）。证明，如果 \(E'\) 与 \(E''\) 都是伪凝聚的（相应地，凝聚的），则 E 亦然。证明，如果 E 与 \(E''\) 是凝聚的，则 \(E'\) 亦然（使用 X，第 11 页，命题 6）。
\item
  设 E 是一个凝聚 A-模，E' 是一个伪凝聚（相应地，凝聚）A-模。证明：对每个同态 \(u: \mathrm{E} \to \mathrm{E}'\) , \(\operatorname{Im}(u)\) 与 \(\operatorname{Ker}(u)\) 都是凝聚的，且 Coker（ \(u\) ）是伪凝聚的（相应地，凝聚的）（使用 \(a\) )）。
\item
  证明：伪凝聚（相应地，凝聚）模的任意直和（相应地，任意有限直和）是伪凝聚（相应地，凝聚）模。
\item
  若 E 是伪凝聚模，且 M, N 是 E 的凝聚子模，证明 \(M + N\) 与 \(M \cap N\) 是凝聚的（使用 a) 和 c)）。
\item
  设 A 是交换的。证明：若 E 是凝聚 A-模，F 是凝聚（相应地，伪凝聚）A-模，则 Hom\_A(E, F) 是凝聚（相应地，伪凝聚）A-模。（化归到 F 是凝聚的情形；考虑 E 的一个有限表示，然后使用 b)。）
\end{enumerate}

¶11) a) 设 A 是一个环。证明下列性质是等价的：

\(\alpha)\) A-模 \(\mathbf{A}_s\) 是凝聚的（习题 10）。

β) 每个有限表示的 A-模都是凝聚的。

γ) 对每个有限表示的右 A-模 M，对偶 M* 是有限生成的 A-模。

δ) 对每个集合 I，右 A-模 \(A_{d}^{I}\) 是平坦的。

ε) 平坦右 A-模的任意乘积是平坦的。

（为了证明 \(\alpha\) ）蕴含 \(\beta\) ），使用习题 10，b)。为了证明 \(\delta\) ）蕴含 \(\gamma\) ），使用 X，第 14 页，定理 1；为了证明 \(\alpha\) ）蕴含 \(\varepsilon\) ），使用 X，第 170 页，习题 18。）

满足这些性质的环称为左凝聚的（或简称为凝聚的），右凝聚环的概念类似地定义。

\begin{enumerate}
\def\labelenumi{\alph{enumi})}
\setcounter{enumi}{1}
\item
  证明左诺特环是凝聚的。给出一个不是凝聚环的右阿廷环的例子（参见 VIII，§ 1，习题 4）。
\item
  证明绝对平坦环（X，第 169 页，习题 16）是凝聚的。
\item
  证明，若 A 是凝聚环且 M 是 A-模，则我们有 \(\lambda(M) = -1\) 或 \(\lambda(M) = 0\) 或 \(\lambda(M) = +\infty\) （习题 6）。特别地，每个有限表示的 A-模都容许一个由有限生成自由模构成的分解。
\item
  设 \((\mathbf{A}_{\alpha}, \varphi_{\beta \alpha})\) 为一个环的归纳系统，其指标集是滤过的，并设 \(\mathbf{A} = \varinjlim \mathbf{A}_{\alpha}\) 。我们假设对 \(\alpha \leqslant \beta\) , \(\mathbf{A}_{\beta}\) 是一个平坦的右 \(\mathbf{A}_{\alpha}\) -模。证明若诸 \(\mathbf{A}_{\alpha}\) 是凝聚的，则 \(\mathbf{A}\) 亦然（注意 \(\mathbf{A}\) 是一个平坦的右 \(\mathbf{A}_{\alpha}\) -模，对所有 \(\alpha\) 成立，且对每个有限生成的左理想 \(\alpha \subset \mathbf{A}\) ，存在一个指标 \(\alpha\) 与一个理想 \(\alpha_{\alpha}\) （ \(\mathbf{A}\) 中的），使得 \(\alpha\) 同构于 \(\mathbf{A} \otimes_{\mathbf{A}} a_{\alpha'}\) 。）
\item
  从 e) 推出，交换诺特环上的任意多项式环（对任意有限或无限个未定元）是凝聚的。由此推出，凝聚环的商环不一定是凝聚的。
\item
  为使 A 是左凝聚的，必要且充分条件是 A 的每个元素的左零化子是有限生成的，并且两个有限生成的左理想的交是有限生成的（使用习题 6，f））。
\end{enumerate}

\begin{enumerate}
\def\labelenumi{\arabic{enumi})}
\setcounter{enumi}{11}
\item
  设 c 为无穷基数。假设环 A 的每个左理想都容许一个基数 ≤ c 的生成集。证明：由基数 ≤ c 的元素集生成的每个 A-模都容许一个由秩 ≤ c 的自由 A-模构成的左分解（参见 VIII，§ 1，习题 13）。
\item
  设 M 是一个 A-模。M 的一个极小内射分解是 M 的一个内射分解 \((e, 1)\) ，使得 \(I^n\) 是 \(Z^n(I)\) 的内射包络（对所有 \(n \geqslant 0\) ）。
\end{enumerate}

\begin{enumerate}
\def\labelenumi{\alph{enumi})}
\tightlist
\item
  每个 A-模都具有极小内射分解；任意两个这样的分解都是同构的。
\item
  设 I、I' 为 M 的两个内射分解，其中 I 是极小的；设 \(f: I \to I'\) 与 \(g: I' \to I\) 为两个分解的态射。则 f 是单射，g 是满射，且它是子复形 Im(f)（同构于 I）与 Ker(g)（带零同调）的直和。
\end{enumerate}

\begin{enumerate}
\def\labelenumi{\arabic{enumi})}
\setcounter{enumi}{13}
\tightlist
\item
  假设环 A 是局部诺特环；用 m 表示其极大理想。设 L 是由有限生成 A-模组成的复形，它在次数 \textless{} 0 时为零，M 是一个有限生成的 A-模，且 q : L → M 是一个态射。我们作如下假设：
\end{enumerate}

\begin{enumerate}
\def\labelenumi{(\roman{enumi})}
\item
  同态 \(1 \otimes q : (\mathbf{A}/\mathfrak{m}) \otimes_{\mathbf{A}} \mathbf{L}_{0} \to (\mathbf{A}/\mathfrak{m}) \otimes_{\mathbf{A}} \mathbf{M}\) 是单射。
\item
  我们有 \(d_{i}(\mathbf{L}_{i})\subset \mathfrak{m}\mathbf{L}_{i - 1}\) 当 \(i\geqslant 1\) ，并且映射
\end{enumerate}

\[
\bar {d} _ {i}: \mathrm{L} _ {i} / \mathrm{mL} _ {i} \rightarrow \mathrm{mL} _ {i - 1} / \mathrm{m} ^ {2} \mathrm{L} _ {i - 1} \quad \text { is injective }.
\]

设 (P, p) 是 M 的一个投射分解， \(u : L \to P\) 为一个态射，使得 \(p \circ u = q\) 。证明 u 是单射，并且 \(u(L)\) 是 P 作为分次 A-模的一个直和项（化归到分解 P 为极小的情形，并利用 VIII，§ 8，第 3 号，推论 3）。

15) 设 A 是一个诺特局部环，m 为其极大理想，k = A/m。设 P 为 A-模 k 的一个极小投射分解；对 \(n \geqslant 0\) ，我们用 \(b_{n}\) 表示自由 A-模 \(P_{n}\) 的秩。

\begin{enumerate}
\def\labelenumi{\alph{enumi})}
\item
  证明我们有 \(b_0 = 1\) ，并且 \(b_1 = \dim_k (m/m^2)\) 是 \(m\) 的生成元的最小数目。
\item
  若 A 是交换的，我们有 \(b_{2} \geqslant \frac{1}{2} b_{1}(b_{1} - 1)\) （考虑正合列 \(\mathbf{P}_{2} \to m\mathbf{P}_{1} \to m^{2} \to 0\) ）。 c) 假设环 A 是阿廷环（不一定是交换的）。证明不等式 \(b_{2} \geqslant \frac{1}{4} b_{1}^{2}\) （设 \(\mathbf{M}_{l}\) 为 \(\mathbf{P}_{2}\) 的由元素 \(x\) （满足 \(d_{2}(x) \in m^{l}\mathbf{P}_{1}\) ）组成的子模；证明我们有
\end{enumerate}

\[
\mathbf {P} _ {2} \subset \mathbf {M} _ {l + 1}
\]

并且存在一个正合列 \(0 \to \mathbf{M}_l / \mathbf{M}_{l+1} \to \mathfrak{m}^l \mathbf{P}_1 / \mathfrak{m}^{l+1} \mathbf{P}_1 \to \mathfrak{m}^{l+1} / \mathfrak{m}^{l+2} \to 0\) 。由此推出多项式 \(p(t) = \sum_{l} (\dim_k (\mathfrak{m}^l / \mathfrak{m}^{l+1})) t^l\) 满足 \((b_2 t^2 - b_1 t + 1)p(t) \geqslant 0\) （当 \(0 \leqslant t \leqslant 1\) ）。通过分别考察情形

 \(b_{1} = 1, b_{1} \geqslant 2\) ，得出结论。

\begin{enumerate}
\def\labelenumi{\arabic{enumi})}
\setcounter{enumi}{15}
\tightlist
\item
  \begin{enumerate}
  \def\labelenumii{\alph{enumii})}
  \tightlist
  \item
    设 \(0 \to \mathbf{M}' \to \mathbf{M} \to \mathbf{M}'' \to 0\) 为 A-模的一个正合列， \((\mathbf{P}', p')\) （相应地， \((\mathbf{P}'', p'')\) ）为 \(\mathbf{M}'\) （相应地， \(\mathbf{M}''\) ）的一个投射分解。证明存在一个投射分解 \((\mathbf{P}, p)\) （ \(\mathbf{M}\) 的）以及一个交换图：
  \end{enumerate}
\end{enumerate}

\[
\begin{array}{c} 0 \to \mathrm{P} ^ {\prime} \to \mathrm{P} \to \mathrm{P} ^ {\prime \prime} \to 0 \\ \downarrow^ {p ^ {\prime}} \quad \downarrow^ {p} \quad \downarrow^ {p ^ {\prime \prime}} \\ 0 \to \mathrm{M} ^ {\prime} \to \mathrm{M} \to \mathrm{M} ^ {\prime \prime} \to 0 \end{array}
\]

其中水平行是正合的。

\begin{enumerate}
\def\labelenumi{\alph{enumi})}
\setcounter{enumi}{1}
\tightlist
\item
  设 \(0 \to \mathbf{N}' \to \mathbf{N} \to \mathbf{N}'' \to 0\) 为 A-模的第二个正合列， \((\mathbf{Q}, q)\) （相应地， \((\mathbf{Q}', q')\) ，相应地， \((\mathbf{Q}'', q'')\) ）为 \(\mathbf{N}\) （相应地， \(\mathbf{N}'\) ，相应地， \(\mathbf{N}''\) ）的一个投射分解，从而存在一个具有正合行的交换图：
\end{enumerate}

\[
\begin{array}{c}0 \rightarrow Q ^ {\prime} \rightarrow Q \rightarrow Q ^ {\prime \prime} \rightarrow 0\\\downarrow^ {q ^ {\prime}} \quad \downarrow^ {q} \quad \downarrow^ {q ^ {\prime \prime}}\\0 \rightarrow N ^ {\prime} \rightarrow N \rightarrow N ^ {\prime \prime} \rightarrow 0\end{array}
\]

另一方面，设：

\[
\begin{array}{c} 0 \to \mathrm{M} ^ {\prime} \to \mathrm{M} \to \mathrm{M} ^ {\prime \prime} \to 0 \\ \downarrow^ {u ^ {\prime}} \quad \downarrow^ {u} \quad \downarrow^ {u ^ {\prime \prime}} \\ 0 \to \mathrm{N} ^ {\prime} \to \mathrm{N} \to \mathrm{N} ^ {\prime \prime} \to 0. \end{array}
\]

它是正合列的一个交换图；设 \(\tilde{u}':\mathbf{P}'\to\mathbf{Q}'\) 与 \(\tilde{u}'':\mathbf{P}''\to\mathbf{Q}''\) 为复形的态射，使得 \(u'\circ p'=q'\circ\tilde{u}'\) 与 \(u''\circ p''=q''\circ\tilde{u}'\) 。证明存在一个态射 \(\tilde{u}:\mathbf{P}\to\mathbf{Q}\) ，使得 \(u\circ p=q\circ\tilde{u}\) ，并且下图：

\[
0 \rightarrow \mathrm{P} ^ {\prime} \rightarrow \mathrm{P} \rightarrow \mathrm{P} ^ {\prime \prime} \rightarrow 0\tag{*}
\]

\[
0 \rightarrow Q ^ {\prime} \rightarrow Q \rightarrow Q ^ {\prime \prime} \rightarrow 0
\]

是交换的。

\begin{enumerate}
\def\labelenumi{\alph{enumi})}
\setcounter{enumi}{2}
\tightlist
\item
  设 \(\tilde{u}_{1}^{\prime}: \mathrm{P}^{\prime} \to \mathrm{Q}^{\prime}, \tilde{u}_{1}: \mathrm{P} \to \mathrm{Q}\) 与 \(\tilde{u}_{1}^{\prime\prime}: \mathrm{P}^{\prime\prime} \to \mathrm{Q}^{\prime\prime}\) 为三个复形的态射，使得
\end{enumerate}

\[
u ^ {\prime} \circ p ^ {\prime} = q ^ {\prime} \circ \tilde {u} _ {1} ^ {\prime}, u \circ p = q \circ \tilde {u} _ {1}, u ^ {\prime \prime} \circ p ^ {\prime \prime} = q ^ {\prime \prime} \circ \tilde {u} _ {1} ^ {\prime \prime}
\]

使图 (*) 交换；设 \(s'\) （相应地， \(s''\) ）为连接 \(\tilde{u}'\) 与 \(\tilde{u}_{1}'\) （相应地， \(\tilde{u}''\) 与 \(\tilde{u}_{1}'\) ）的一个同伦。证明存在一个同伦 \(s\) 连接 \(\tilde{u}\) 与 \(\tilde{u}_{1}\) ，使得该图：

\[
\begin{array}{c}0 \rightarrow \mathrm{P} ^ {\prime} \rightarrow \mathrm{P} \rightarrow \mathrm{P} ^ {\prime \prime} \rightarrow 0\\\downarrow s ^ {\prime} \quad \downarrow s \quad \downarrow s ^ {\prime \prime}\\0 \rightarrow \mathrm{Q} ^ {\prime} \rightarrow \mathrm{Q} \rightarrow \mathrm{Q} ^ {\prime \prime} \rightarrow 0\end{array}
\]

是交换的。

\begin{enumerate}
\def\labelenumi{\alph{enumi})}
\setcounter{enumi}{3}
\tightlist
\item
  陈述并证明关于内射分解的与 a)、b)、c) 相对应的结果。 ¶17) 在本习题中，若 \((\mathbf{K}, d', d'')\) 是一个双复形（X，第 174 页，习题 11），我们将用 \(K_{p,*}\) 表示复形 \((\bigoplus_{n} K_{p,n}, d'')\) 。每个复形都将被视为一个双复形，它在次数 \((p, q)\) （当 \(q \neq 0\) ）处为零。
\end{enumerate}

设 C 为 A-模的一个复形。C 的一个投射分解是一个二元组 \((P, p)\) ，其中 P 是一个双复形且 \(p : P \to C\) 是双复形的一个态射，使得 \(P_{p,q} = 0\) （对 q \textless{} 0）且使得复形 \(P_{p,*}\) （相应地， \(Z'_{p,*}(\hat{P})\) , \(B'_{p,*}(\mathbf{P})\) , \(H'_{p,*}(\mathbf{P})\) ）定义了 \(C_p\) （相应地， \(Z_p(\mathbf{C})\) , \(B_p(\mathbf{C})\) , \(H_p(\mathbf{C})\) ）的一个投射分解（对每个 \(p \in Z\) ）。为使 \((P, p)\) 是 C 的一个投射分解，只需复形 \(B'_{p,*}(\mathbf{P})\) 与 \(H'_{p,*}(\mathbf{P})\) 对每个 p 定义 \(B_p(\mathbf{C})\) 与 \(H_p(\mathbf{C})\) 的投射分解即可。

\begin{enumerate}
\def\labelenumi{\alph{enumi})}
\tightlist
\item
  证明每个 A-复形 C 都容许一个投射分解（对每个 p，选取 \(B_{p,*}\) 与 \(H_{p,*}\) 为 \(B_{p}(C)\) 与 \(H_{p}(C)\) 的投射分解；利用习题 16，构造投射分解 \(Z_{p,*}\) （ \(Z_{p}(C)\) 的）与 \(P_{p,*}\) （ \(C_{p}\) 的），以及正合列：
\end{enumerate}

\[
0 \rightarrow \mathrm{B} _ {p, *} \rightarrow \mathrm{Z} _ {p, *} \rightarrow \mathrm{H} _ {p, *} \rightarrow 0; \quad 0 \rightarrow \mathrm{Z} _ {p, *} \rightarrow \mathrm{P} _ {p, *} \rightarrow \mathrm{B} _ {p + 1, *} \rightarrow 0).
\]

\begin{enumerate}
\def\labelenumi{\alph{enumi})}
\setcounter{enumi}{1}
\item
  设 \(\mathbf{C}'\) 为另一个复形， \((\mathbf{P}', p')\) 为 \(\mathbf{C}'\) 的一个投射分解， \(u: \mathbf{C} \to \mathbf{C}'\) 为复形的一个态射。证明存在一个双复形的态射 \(\tilde{u}: \mathbf{P} \to \mathbf{P}'\) ，使得 \(u \circ p = p' \circ \tilde{u}\) （对每个 \(p \in \mathbf{Z}\) ，选取态射 \(\mathbf{B}_{p'}', (\mathbf{P}) \to \mathbf{B}_{p'}', (\mathbf{P}')\) 与 \(\mathrm{H}_{p'}', (\mathbf{P}) \to \mathrm{H}_{p'}', (\mathbf{P}')\) ；如前所述构造 \(\tilde{u}\) ，使用习题 16，b)）。
\item
  设 \(v: \mathbf{C} \to \mathbf{C}'\) 为与 \(u\) 同伦的复形的态射，并设 \(\tilde{v}: \mathbf{P} \to \mathbf{P}'\) 为一个双复形的态射，使得 \(v \circ p = p' \circ \tilde{v}\) 。证明 \(\tilde{u}\) 与 \(\tilde{v}\) 是同伦的（按 X，第 174 页，习题 11 c) 的意义；利用习题 16 c)）。
\end{enumerate}

\(d)\) 设 E 为 \(\mathbf{Z}\) 的一个子集，且 C 为一个 A-复形，使得 \(\mathbf{C}_{p}=0\) （当 \(p\in\mathbf{E}\) ）。证明存在 C 的一个投射分解 \((\mathbf{P},p)\) ，使得 \(\mathbf{P}_{p,q}=0\) （当 \(p\in\mathbf{E}\) ）。

\begin{enumerate}
\def\labelenumi{\alph{enumi})}
\setcounter{enumi}{4}
\item
  假设存在一个整数 \(n\) ，使得每个 A-模都容许一个长度为 \(\leqslant n\) 的投射分解。证明每个 A-复形 C 都容许一个投射分解 \((P, p)\) ，使得 \(P_{p,q} = 0\) （当 \(q > n\) ）。
\item
  设 C 是一个 A-复形。C 的一个内射分解是一个二元组 \((e, l)\) ，其中 I 是一个双复形且 \(e: C \to 1\) 是双复形的一个态射，使得 \(I^{p,q} = 0\) （对 q \textless{} 0）且使得复形 \(I^{p,*}\) （相应地， \(Z^{p,*}(I)\) , \(B^{p,*}(l)\) , \(H^{p,*}(l)\) ）定义了 \(C^{p}\) （相应地， \(Z^{p}(C)\) , \(B^{p}(C)\) , \(H^{p}(C)\) ）的一个内射分解（对每个 \(p \in Z\) ）。陈述并证明内射分解情形下结果 a) 到 e) 的类比。
\end{enumerate}

\begin{enumerate}
\def\labelenumi{\arabic{enumi})}
\setcounter{enumi}{17}
\tightlist
\item
  假设环 A 是交换的。一个微分分次 A-代数是指一个分次类型为 Z 的 A-代数 S，使得 \(S_{n} = 0\) （对 n \textless{} 0），并配备一个次数为 (-1) 的反导子 d，满足 \(d \circ d = 0\) 。我们也用 S 表示 A-模复形 (S, d)。
\end{enumerate}

\begin{enumerate}
\def\labelenumi{\alph{enumi})}
\item
  设 C 为 A-模复形，在次数 \textless{} 0 处为零。证明在 S = T(C)（相应地 S = S(C)，相应地 S = Λ(C)）上存在一个微分分次 A-代数结构，使得 C 到 S 的典范嵌入是复形的态射。
\item
  证明 S 的乘法在 H(S) 上诱导一个分次 A-代数结构。
\item
  设 S、T 为两个微分分次 A-代数。令 D(s ⊗ t) = ds ⊗ t + (-1)\^{}p s ⊗ dt（对 s ∈ S\_p, t ∈ T）。证明 D 在左张量积 S *⊗ \_A T（III，第 49 页）上定义了微分分次 A-代数的结构。
\end{enumerate}

\begin{enumerate}
\def\labelenumi{\arabic{enumi})}
\setcounter{enumi}{18}
\tightlist
\item
  假设环 A 是交换的。
\end{enumerate}

\begin{enumerate}
\def\labelenumi{\alph{enumi})}
\item
  设 \((\mathbb{S}, d_{\mathbb{S}})\) 为一个微分分次 A-代数（习题 18），M 是一个 A-模， \(u: \mathbb{M} \to Z_{n}(\mathbb{S})\) 为一个 A-同态。证明在 \(\mathbb{S} \otimes_{\mathbb{A}} \mathbf{T}(\mathbb{M})\) 上存在一个微分分次 A-代数结构，使得 \(d(s \otimes 1) = d_{\mathbb{S}} s \otimes 1\) （当 \(s \in \mathbb{S}\) ）且 \(d(1 \otimes m) = u(m) \otimes 1\) （当 \(m \in \mathbb{M}\) ）。
\item
  设 B 是一个 A-代数。证明存在一个微分分次 A-代数 S 和一个 A-代数同态 \(p: \mathbf{S}_{0} \to \mathbf{B}\) ，使得 \((\mathbf{S}, p)\) 是 A-模 B 的一个自由分解。（通过对 \(n\) 作归纳，利用 \(a\) )，构造一个微分分次代数 \(\mathbf{S}(n)\) ，它在每个次数都是自由的，使得 \(\mathrm{H}_{i}(\mathbf{S}(n)) = 0\) （当 \(0 < i < n\) ）且 \(\mathrm{H}_{0}(\mathbf{S}(n)) = \mathbf{B}\) 。）
\item
  若代数 B 是交换的，证明可以找到一个交错 A-代数 S（III，第 53 页）满足 b) 的条件。若进一步 A 是诺特的，且 B 是 A 关于一个理想的商，证明可以找到 S 使得 \(S_{0} = A\) 并且 \(S_{n}\) 对每个 n 都是有限生成的 A-模。
\end{enumerate}

¶ 20) 假设环 A 是交换环 k 上的一个代数。设 \(\eta: k \to A\) 为由 \(\eta(\lambda) = \lambda 1_{A}\) （当 \(\lambda \in k\) ）定义的同态，并设 \(\overline{A}\) 为其余核。若 \(a\) 是 A 的一个元素，我们用 \(\overline{a}\) 表示它在 \(\overline{A}\) 中的类。我们设 \(\mathrm{N}_{n}(\mathrm{A}) = \mathrm{A} \otimes_{k} \overline{\mathrm{A}}^{\otimes n} \otimes_{k} \mathrm{A}\) （当 \(n \geqslant 0\) ）, \(\mathrm{N}_{n}(\mathrm{A}) = 0\) （当 \(n < 0\) ），且

\[
\mathrm{N} (\mathbf {A}) = \bigoplus_ {n \in \mathbf {Z}} \mathrm{N} _ {n} (\mathbf {A}).
\]

\begin{enumerate}
\def\labelenumi{\alph{enumi})}
\tightlist
\item
  证明存在 N(A) 的一个分次次数为 (-1) 的 (A, A)-自同态 d，使得
\end{enumerate}

\[
\begin{array}{l} d (a \otimes \bar {a} _ {1} \otimes \dots \otimes \bar {a} _ {n} \otimes b) = a a _ {1} \otimes \bar {a} _ {2} \otimes \dots \otimes \bar {a} _ {n} \otimes b + \sum_ {i = 1} ^ {n - 1} (- 1) ^ {i} a \otimes \dots \otimes \overline {{a _ {i}}} \overline {{a _ {i + 1}}} \otimes \dots \otimes \bar {a} _ {n} \otimes b \\ \qquad + (- 1) ^ {n} a \otimes \bar {a} _ {1} \otimes \dots \otimes \bar {a} _ {n - 1} \otimes a _ {n} b \quad \text {for} \quad a, a _ {1},..., a _ {n}, b \in A \end{array}
\]

\begin{enumerate}
\def\labelenumi{\alph{enumi})}
\setcounter{enumi}{1}
\item
  证明 \(d \circ d = 0\) ，使得 \(\mathbf{N}(\mathbf{A})\) 是 \((\mathbf{A}, \mathbf{A})\) -双模的一个复形。若 M 是一个左 A-模，我们记 \(\mathbf{N}(\mathbf{A}, \mathbf{M})\) 为 A-复形 \(\mathbf{N}(\mathbf{A}) \otimes_{\mathbf{A}} \mathbf{M}\) 。
\item
  设 B(A, M) 为 M 的标准分解（X，第 58 页）， \(p: \mathrm{B}(\mathrm{A}, \mathrm{M}) \to \mathrm{N}(\mathrm{A}, \mathrm{M})\) 为由 \(p(a_0 \otimes \cdots \otimes a_n \otimes m) = a_0 \otimes \overline{a}_1 \otimes \cdots \otimes \overline{a}_n \otimes m\) （当 \(a_0, \ldots, a_n \in \mathbf{A}, m \in \mathbf{M}\) ）定义的 A-同态。证明 \(p\) 是 A-复形的一个同伦等价（通过归纳构造一个在同伦意义下的逆映射）。
\end{enumerate}

* 21) 设 \(k\) 是一个交换环，并设 g 是一个 \(k\) 李代数，U 是其包络代数（参见 LIE, I, § 2）。假设 g 是一个 \(k\) -自由模。

\begin{enumerate}
\def\labelenumi{\alph{enumi})}
\tightlist
\item
  证明对每个 \(n \geqslant 0\) 存在一个单射 U-同态
\end{enumerate}

\[
j _ {n}: \mathrm{U} \otimes_ {k} \Lambda^ {n} (\mathrm{g}) \rightarrow \mathrm{U} ^ {\otimes (n + 1)} \quad \text { such that } \quad j _ {n} (u \otimes (x _ {1} \wedge \dots \wedge x _ {n})) = \sum_ {\sigma \in \mathfrak {S} _ {n}} \varepsilon (\sigma), u \otimes x _ {\sigma (1)} \otimes \dots \otimes x _ {\sigma (n)}
\]

当 \(u\in \mathbf{U},x_1,\dots ,x_n\in \mathfrak{g}\)

因此我们定义一个次数为 0 的分次 U-同态 \(j:U\otimes_{k}\Lambda(g)\to B(U,k)\) ，其中 \(B(U,k)\) 是左 U-模 \(k\) 的标准分解（ \(k\) 的 U-模结构由自然同态 \(U\to k\) 导出）。证明 \(U\otimes_{k}\Lambda(g)\) 通过 \(j\) 等同于 \(B(U,k)\) 的一个子复形，其微分由下述公式给出：

\[
\begin{array}{l} d (u \otimes (x _ {1} \wedge \dots \wedge x _ {n})) = \sum_ {1 \leqslant i \leqslant n} (- 1) ^ {i + 1} u x _ {i} \otimes (x _ {1} \wedge \dots \wedge \hat {x} _ {i} \wedge \dots \wedge x _ {n}) \\ + \sum_ {1 \leqslant i <   j \leqslant n} (- 1) ^ {i + j} u \otimes ([ x _ {i}, x _ {j} ] \wedge x _ {1} \wedge \dots \wedge \hat {x} _ {i} \wedge \dots \wedge \hat {x} _ {j} \wedge \dots \wedge x _ {n}) \quad \text { pour } \quad u \in U, x _ {1},..., x _ {n} \in g. \end{array}
\]

（字母上方的符号表示该字母应被省略）。我们用 V(g) 表示如此定义的 U-复形。b) 证明 V(g) 定义了 U-模 k 的一个自由分解（令 \((\mathrm{U}_{n})_{n\geqslant0}\) 为 U 的自然滤过， \(F_{r}V(g)\) 为 V(g) 的由元素 \(u_{p}\otimes x_{q}\) （当 \(u_{p}\in U_{p}, x_{q}\in\Lambda^{q}(g), p+q\leqslant r\) ）生成的子 k-模。证明 \(F_{r}V(g)\) 是 V(g) 的一个子复形，并且复形 \(\bigoplus_{r\geqslant0}F_{r}V(g)/F_{r-1}V(g)\)

同构于复形 \(\mathbf{S}(\mathfrak{g})\otimes\symbf{\Lambda}(\mathfrak{g})\) （在 X，第 151 页中定义）。*
% semantic-fragments: legacy-input-seam
\relax{}
